% Bagean: metodhe
% Bab: bukti
% Pérangan: migunakaké dhéfinisi

\documentclass[../../../include/open-logic-section]{subfiles}

\begin{document}

\olfileid{mth}{prf}{def}

\olsection{Migunakaké Dhéfinisi}

Kita wis nyebut yèn panjenengan kudu ngerti kabèh dhéfinisi kang
bisa digunakaké ing bukti lan bisa nrapaké kanthi bener. Bab iki
wigati banget, mula prayoga ditliti luwih rinci. Dhéfinisi
nyingket sipat lan relasi supaya luwih gampang dirembug. Ungkapan
cekak kang diweruhaké diarani \emph{definiendum}, déné bab kang
disingket diarani \emph{definiens}. Ing bukti, kita asring kudu
mbalèkaké definiendum marang dhéfinisi asalé, amarga struktur
logis saka definiens---wujud jangkep kang disingket---iku kang
mbantu kita ngrampungaké bukti. Kanthi njlentrehaké dhéfinisi,
panjenengan tekan pérangan kang ngandhut panalaran logis wigati.

Kita miwiti kanthi conto. Anggepen panjenengan arep mbuktekaké:

\begin{prop}
Kanggo sembarang himpunan $A$ lan $B$, $A \cup B = B \cup A$.
\end{prop}

Kanggo miwiti bukti waé, kita kudu ngerti tegesé rong himpunan
padha, yaiku tegesé ``$=$'' ing persamaan iku kanggo himpunan.
Rong himpunan diarani padha yèn nduwèni !!{element}s kang padha.
Dadi dhéfinisi kang
kudu kita njlentrehaké yaiku:

\begin{defn}
Himpunan $A$ lan $B$ \emph{padha}, $A = B$, yèn lan mung yèn
saben !!{element} saka~$A$ uga !!a{element} saka~$B$, lan kosok baliné.
\end{defn}

Dhéfinisi iki migunakaké $A$ lan~$B$ minangka panggonan kanggo
sebarang himpunan. Bab kang didhéfinisikaké---\emph{definiendum}---
yaiku ungkapan ``$A = B$'', kanthi mènèhi syarat kapan $A = B$
bener. Syarat iki---``saben !!{element} saka~$A$ iku
!!a{element} saka~$B$, lan kosok baliné''---yaiku
\emph{definiens}.\footnote{Ing kasus iki, kang bisa mbingungaké,
yèn $A = B$, himpunan $A$ lan $B$ sakjané sawijining himpunan
kang padha, sanadyan aksarané béda ing sisih kiwa lan tengen.
Nanging carané ngenali himpunan mau bisa béda, mula dhéfinisi
iki tetep wigati.} Dhéfinisi iku negesaké yèn $A = B$ bener yèn
lan mung yèn syarat mau kelakon; iki sok disingket dadi ``iff''.

Nalika nrapaké dhéfinisi, panjenengan kudu nyocogaké $A$ lan $B$
ing dhéfinisi karo kasus kang lagi dirembug. Ing conto iki,
supaya $A \cup B = B \cup A$ bener, saben $z \in A \cup B$
uga kudu ana ing $B \cup A$, lan kosok baliné. Ungkapan $A \cup
B$ ing proposisi dadi pengganti~$A$ ing dhéfinisi, déné
$B \cup A$ dadi pengganti~$B$. Amarga $A$ lan $B$ padha dienggo
ing dhéfinisi lan ing proposisi nanging kanthi peran béda,
panjenengan kudu ngati-ati supaya ora kecampur. Upamané,
mbuktekaké yèn saben !!{element} saka~$A$ iku !!a{element}
saka~$B$ lan kosok baliné ora mbuktekaké proposisi iki;
iku mung mbuktekaké $A = B$, dudu $A \cup B = B \cup A$.
Kajaba iku, amarga $A$ lan $B$ bisa sembarang himpunan,
tanpa anggepan tambahan $A$ lan $B$ bisa waé béda.

Ing bukti iki kita uga migunakaké pangertèn teori himpunan kaya
gabungan, mula kudu ngerti tegesé simbol $\cup$ supaya ngerti
carané nerusaké bukti. Kadhangkala njlentrehaké sawijining
dhéfinisi mbutuhaké njlentrehaké dhéfinisi liya maneh. Upamané,
$A \cup B$ didhéfinisikaké minangka $\Setabs{z}{z \in A
  \text{ utawa } z \in B}$. Yèn arep mbuktekaké $x \in A \cup
B$, dhéfinisi $\cup$ nuduhaké yèn kudu mbuktekaké
$x \in \Setabs{z}{z \in A \text{ utawa } z \in B}$.
Saiki élinga uga yèn $x \in \Setabs{z}{\dots z\dots}$ yèn lan
mung yèn $\dots x\dots$. Dadi sawisé notasi
$\Setabs{z}{\dots z \dots}$ uga dijlentrehaké, kang kudu
dibuktèkaké yaiku $x \in A$ utawa $x \in B$. Mula pratelan
``saben !!{element} saka $A \cup B$ uga !!a{element} saka
$B \cup A$'' tegesé ``kanggo saben $x$, yèn $x \in A$
utawa $x \in B$, mula $x \in B$ utawa $x \in A$''.
Yèn kabèh dhéfinisi ing proposisi dijlentrehaké, kang kudu
dibuktèkaké dadi:

\begin{prop}
Kanggo sembarang himpunan $A$ lan $B$: (a) kanggo saben $x$,
yèn $x \in A$ utawa $x \in B$, mula $x \in B$ utawa $x \in A$;
lan (b) kanggo saben $x$, yèn $x \in B$ utawa $x \in A$,
mula $x \in A$ utawa $x \in B$.
\end{prop}

Bab kang wigati yaiku njlentrehaké dhéfinisi minangka pérangan
kang perlu nalika nyusun bukti. Nindakaké iki kanthi bener
kadhangkala angel: variabel ing dhéfinisi kudu dibédakaké lan
dicocogaké karo suku ing pratelan kang dibuktèkaké. Supaya
kasil, panjenengan kudu ngerti apa kang dijaluk pitakonan lan
tegesé saben istilah ing kono; asring luwih saka siji dhéfinisi
kudu dijlentrehaké. Ing bukti prasaja kaya conto ing ngisor iki,
solusiné meh langsung katut saka dhéfinisi. Nanging ora saben
bukti segampang iku.

\begin{prob}
Anggepen panjenengan dijaluk mbuktekaké $A \cap B \neq \emptyset$.
Jlentrehna kabèh dhéfinisi ing kéné: pratelakna maneh tanpa
nyebut ``$\cap$'', ``='', utawa ``$\emptyset$''.
\end{prob}

\end{document}
